
% ==================== ВВЕДЕНИЕ ====================
\section*{Введения}
\addcontentsline{toc}{section}{Введения}

В современных условиях экспоненциального роста количества информации в таких областях как здравоохранение, фармакология и биомедицина, где ежедневно пишутся десятки тысяч статей и патентов, данные остаются неструктурированными. В узкоспециализированных предметных областях, таких как хронобиология, ручной поиск и мониторинг научных публикаций и выявление неочевидных связей между веществами становится невозможным из-за большого объёма информации.

На данный момент существует единственная в мире структурированная база данных по хронобиотикам ChronobioticsDB, созданная для структуризации работы в данной области. Существует веб-интерфейс этой базы данных, однако аналитический потенциал ограничен, так как это реляционная база данных, работающая с SQL.

Предмет исследования --- методы и технологии автоматизации анализа научных данных с использованием искусственного интеллекта.

Практическая значимость проекта заключается в создании удобного и функционального инструмента, который обеспечит эффективный доступ к уникальной базе данных хронобиотиков.

\newpage

% ==================== ГЛАВА 1 ====================
\section{Теоретико-методологические основы построения агентных систем для анализа научных данных}

\subsection{Эволюция парадигм интеллектуальных систем от экспертных систем к мультиагентным архитектурам}

\subsubsection{Исторический анализ развития экспертных систем в научных исследованиях}

Эволюция искусственного интеллекта (ИИ) представляет собой сложный путь от попыток формализации логических рассуждений до создания систем, способных к автономным действиям. Двумя ключевыми вехами на этом пути стали экспертные системы, доминировавшие в 1970–1980-х годах, и современный агентно-ориентированный подход (Agentic AI). Анализ их развития позволяет не только проследить смену технологических парадигм, но и выявить фундаментальные принципы, которые остаются неизменными при создании интеллектуальных систем, способных решать сложные, слабо формализуемые задачи.

\subsubsection{Становление агентно-ориентированного подхода в искусственном интеллекте}

В данном анализе рассматривается исторический путь от первых экспертных систем, основанных на статичных правилах, к агентным системам, демонстрирующим целенаправленное рассуждение и автономное выполнение действий.

\subsubsection{Таксономия интеллектуальных агентов по функциональному назначению и архитектурным принципам}

Текст раздела...

\subsubsection{Современные программные фреймворки для конструирования мультиагентных систем}

Текст раздела...

\subsection{Большие языковые модели как фундаментальный компонент интеллектуальных систем}

\subsubsection{Архитектурные принципы и эволюция больших языковых моделей}

Текст раздела...

\subsubsection{Классификация и сравнительный анализ современных LLM для решения научных задач}

Текст раздела...

\subsubsection{Проприетарные модели: архитектура и особенности применения (OpenAI GPT, Anthropic Claude, Google Gemini, Cohere)}

Текст раздела...

\subsubsection{Открытые и исследовательские модели: особенности и преимущества (BLOOM, LLaMA, Mistral, Qwen)}

Текст раздела...

\subsubsection{Модели для локального развертывания: оптимизация и инференс (Local LLM, Custom Model)}

Текст раздела...

\subsubsection{Методология дообучения языковых моделей на предметных данных}

Текст раздела...

\subsubsection{Сравнительный анализ эффективности различных LLM в научных исследованиях}

Текст раздела...

\subsection{Теоретические основы извлечения информации из гетерогенных источников}

\subsubsection{Методы семантического поиска и векторного представления текстов}

Текст раздела...

\subsubsection{Архитектура RAG: принципы и компоненты}

Текст раздела...

\subsubsection{Графовые модели представления знаний в биоинформатике}

Текст раздела...

\subsubsection{Гибридные стратегии поиска информации в реляционных базах данных}

Текст раздела...

\subsection{Технологии мультимодального человеко-машинного взаимодействия}

\subsubsection{Современные методы автоматического распознавания речи}

Текст раздела...

\subsubsection{Технологии синтеза речи и их применение в научных системах}

Текст раздела...

\subsubsection{Интеграция текстовых, голосовых и графических интерфейсов}

Текст раздела...

\subsubsection{Методы обработки голосовых команд в специализированных областях}

Текст раздела...

\subsection{Мультимодальные диалоговые системы: архитектура и принципы}

\subsubsection{Концепция мультимодального диалога}

Текст раздела...

\subsubsection{Методы синхронизации различных модальностей ввода/вывода}

Текст раздела...

\subsubsection{Управление контекстом в мультимодальных системах}

Текст раздела...

\subsubsection{Оценка качества пользовательского опыта в мультимодальных интерфейсах}

Текст раздела...

\subsection{Лингвистическое обеспечение многоязычных интеллектуальных систем}

\subsubsection{Методы автоматической идентификации языка текстов}

Текст раздела...

\subsubsection{Технологии машинного перевода научной терминологии}

Текст раздела...

\subsubsection{Многоязычные языковые модели и их применение}

Текст раздела...

\subsubsection{Геолокационные сервисы и адаптация контента}

Текст раздела...

\subsubsection{Многоязычная поддержка в голосовых интерфейсах}

Текст раздела...

\subsection{Методы компьютерного анализа химических данных}

\subsubsection{Формальные языки описания химических структур}

Текст раздела...

\subsubsection{Методы расчета молекулярных дескрипторов}

Текст раздела...

\subsubsection{Компьютерное зрение для распознавания химических структур из изображений}

Текст раздела...

\subsubsection{Интеграция химических инструментов с LLM}

Текст раздела...

\subsubsection{Голосовой ввод химических формул и названий}

Текст раздела...

\subsection{Хронобиология как предметная область интеллектуального анализа}

\subsubsection{Фундаментальные концепции хронобиологии и хронофармакологии}

Текст раздела...

\subsubsection{Уникальность базы данных хронобиотиков как единственного в мире структурированного источника знаний в предметной области}

Текст раздела...

\subsubsection{Анализ структуры и содержания базы данных хронобиотических веществ}

Текст раздела...

\subsubsection{Специфика анализа биологических ритмов и временных зависимостей}

Текст раздела...

\subsubsection{Терминологические особенности предметной области для голосового ввода}

Текст раздела...

\subsubsection{Научная и практическая значимость автоматизации доступа к уникальным данным}

Текст раздела...

\subsection{Архитектурные паттерны проектирования специализированных агентов}

\subsubsection{Принципы функциональной декомпозиции в агентных системах}

Текст раздела...

\subsubsection{Методы координации и коммуникации между агентами}

Текст раздела...

\subsubsection{Шаблоны проектирования реактивных и когнитивных агентов}

Текст раздела...

\subsubsection{Организация мультимодального взаимодействия на уровне агентов}

Текст раздела...

\subsubsection{Интеграция инструментов расширения функциональности}

Текст раздела...

\subsection{Методология поиска и анализа научной литературы}

\subsubsection{Современные методы автоматизированного поиска научных публикаций}

Текст раздела...

\subsubsection{Интеграция с академическими базами данных}

Текст раздела...

\subsubsection{Методы извлечения ключевой информации из научных статей}

Текст раздела...

\subsubsection{Обнаружение новых публикаций и отслеживание научных трендов}

Текст раздела...

\subsection{Методология оценки эффективности интеллектуальных систем}

\subsubsection{Количественные метрики качества извлечения информации}

Текст раздела...

\subsubsection{Методы валидации результатов на ограниченных выборках}

Текст раздела...

\subsubsection{Критерии качества работы специализированных агентов}

Текст раздела...

\subsubsection{Оценка эффективности мультимодального взаимодействия}

Текст раздела...

\subsection{Анализ существующих решений и постановка задачи исследования}

\subsubsection{Критический анализ существующих систем анализа научных данных}

Текст раздела...

\subsubsection{Идентификация ограничений существующих подходов}

Текст раздела...

\subsubsection{Обоснование необходимости создания интеллектуальной системы для единственной в мире базы хронобиотиков}

Текст раздела...

\subsubsection{Формальная постановка задачи разработки мультимодальной агентной системы}

Текст раздела...

\subsection{Выводы и результаты по главе}

Текст раздела...

\newpage

% ==================== ГЛАВА 2 ====================
\section{Проектирование агентной системы Chronobiotic}

\subsection{Анализ требований и спецификация системы}

\subsubsection{Функциональные требования к системе}

Текст раздела...

\subsubsection{Технические требования и ограничения}

Текст раздела...

\subsubsection{Анализ существующей базы данных хронобиотиков}

Текст раздела...

\subsubsection{Требования к мультимодальному пользовательскому интерфейсу}

Текст раздела...

\subsection{Концептуальная архитектура агентной системы}

\subsubsection{Общая структурная схема платформы}

Текст раздела...

\subsubsection{Иерархическая организация специализированных агентов}

Текст раздела...

\subsubsection{Взаимодействие с базой данных хронобиотиков}

Текст раздела...

\subsubsection{Информационная модель и потоки данных}

Текст раздела...

\subsubsection{Архитектура мультимодального чата}

Текст раздела...

\subsection{Проектирование ядра системы менеджмента агентов}

\subsubsection{Архитектура модуля диспетчеризации задач}

Текст раздела...

\subsubsection{Управление жизненным циклом и реестр агентов}

Текст раздела...

\subsubsection{Мониторинг состояний и производительности}

Текст раздела...

\subsubsection{Организация параллельного выполнения задач}

Текст раздела...

\subsection{Проектирование аналитических агентов}

\subsubsection{Функциональная спецификация агента химического анализа}

Текст раздела...

\subsubsection{Проектирование агента оценки эффективности}

Текст раздела...

\subsubsection{Разработка модели агента анализа взаимодействий}

Текст раздела...

\subsubsection{Архитектура агента предсказания свойств}

Текст раздела...

\subsubsection{Проектирование агента поиска структурных аналогов}

Текст раздела...

\subsection{Проектирование хронобиологических агентов}

\subsubsection{Разработка экспертной модели хронобиолога}

Текст раздела...

\subsubsection{Архитектура агента поиска в базе данных}

Текст раздела...

\subsubsection{Проектирование агента работы с экспериментальными данными}

Текст раздела...

\subsubsection{Функциональная спецификация агента анализа литературы}

Текст раздела...

\subsubsection{Проектирование агента исследования механизмов действия}

Текст раздела...

\subsubsection{Архитектура агента комплексного анализа веществ}

Текст раздела...

\subsection{Проектирование вспомогательных агентов для мультимодального диалога}

\subsubsection{Модель мультимодального диалогового агента}

Текст раздела...

\subsubsection{Проектирование агента генерации объяснений}

Текст раздела...

\subsubsection{Архитектура вопросно-ответного агента}

Текст раздела...

\subsubsection{Разработка агента формирования рекомендаций}

Текст раздела...

\subsubsection{Функциональная спецификация агента реферирования}

Текст раздела...

\subsection{Проектирование исследовательских агентов}

\subsubsection{Архитектура агента поиска экспериментальных исследований}

Текст раздела...

\subsubsection{Проектирование агента генерации гипотез}

Текст раздела...

\subsubsection{Разработка агента обзора литературы}

Текст раздела...

\subsubsection{Функциональная спецификация агента исследования механизмов}

Текст раздела...

\subsubsection{Проектирование агента патентного поиска}

Текст раздела...

\subsection{Проектирование голосовых агентов для мультимодального чата}

\subsubsection{Архитектура агента распознавания речи}

Текст раздела...

\subsubsection{Проектирование агента синтеза речи}

Текст раздела...

\subsubsection{Разработка агента голосового интерфейса}

Текст раздела...

\subsubsection{Функциональная спецификация агента обработки аудиоданных}

Текст раздела...

\subsubsection{Проектирование агента мультимодального взаимодействия}

Текст раздела...

\subsubsection{Агент обработки голосовых команд для навигации}

Текст раздела...

\subsection{Проектирование многоязычных агентов}

\subsubsection{Архитектура агента определения языка}

Текст раздела...

\subsubsection{Проектирование агента перевода контента}

Текст раздела...

\subsubsection{Разработка агента многоязычного диалога}

Текст раздела...

\subsubsection{Функциональная спецификация агента локализации интерфейса}

Текст раздела...

\subsubsection{Многоязычная поддержка в голосовых интерфейсах}

Текст раздела...

\subsection{Проектирование цитационных агентов}

\subsubsection{Архитектура агента построения библиографии}

Текст раздела...

\subsubsection{Проектирование агента извлечения цитат}

Текст раздела...

\subsubsection{Разработка агента форматирования ссылок}

Текст раздела...

\subsubsection{Функциональная спецификация агента отслеживания источников}

Текст раздела...

\subsubsection{Проектирование агента валидации источников}

Текст раздела...

\subsection{Проектирование агентов данных}

\subsubsection{Архитектура агента анализа контента}

Текст раздела...

\subsubsection{Проектирование агента хранения данных}

Текст раздела...

\subsubsection{Разработка агента валидации данных}

Текст раздела...

\subsubsection{Функциональная спецификация агента поиска в БД}

Текст раздела...

\subsubsection{Проектирование агента перехода по ссылкам}

Текст раздела...

\subsubsection{Архитектура агента сбора данных из веб-источников}

Текст раздела...

\subsection{Проектирование модуля анализа химических данных}

\subsubsection{Архитектура подсистемы предсказания свойств}

Текст раздела...

\subsubsection{Проектирование интеграции с внешними химическими базами}

Текст раздела...

\subsubsection{Разработка подсистемы обработки химических изображений}

Текст раздела...

\subsubsection{Проектирование парсеров химических форматов}

Текст раздела...

\subsubsection{Поддержка голосового ввода химических названий}

Текст раздела...

\subsection{Проектирование модуля графов знаний}

\subsubsection{Архитектура подсистемы извлечения сущностей}

Текст раздела...

\subsubsection{Проектирование построения графа знаний}

Текст раздела...

\subsubsection{Разработка механизма гибридного поиска}

Текст раздела...

\subsubsection{Функциональная спецификация подсистемы логического вывода}

Текст раздела...

\subsubsection{Проектирование поиска взаимосвязей}

Текст раздела...

\subsection{Проектирование модуля генерации с дополненной выборкой}

\subsubsection{Архитектура подсистемы сегментации текстов}

Текст раздела...

\subsubsection{Проектирование генерации векторных представлений}

Текст раздела...

\subsubsection{Разработка механизмов поиска}

Текст раздела...

\subsubsection{Функциональная спецификация подсистемы ранжирования}

Текст раздела...

\subsubsection{Проектирование векторных хранилищ}

Текст раздела...

\subsection{Проектирование интеграции с большой языковой моделью}

\subsubsection{Архитектура унифицированного взаимодействия с LLM}

Текст раздела...

\subsubsection{Проектирование адаптеров для различных моделей (OpenAI, Anthropic, Google)}

Текст раздела...

\subsubsection{Разработка интеграции с открытыми моделями (BLOOM, LLaMA, Mistral)}

Текст раздела...

\subsubsection{Проектирование подсистемы управления контекстом и промптами}

Текст раздела...

\subsubsection{Разработка механизмов роутинга запросов между моделями}

Текст раздела...

\subsubsection{Функциональная спецификация оптимизации запросов и кэширования}

Текст раздела...

\subsubsection{Проектирование системы выбора оптимальной модели для задачи}

Текст раздела...

\subsubsection{Разработка инструментов расширения функциональности (LLM Tools)}

Текст раздела...

\subsection{Проектирование модуля мультимодального чата}

\subsubsection{Архитектура чат-движка с поддержкой нескольких модальностей}

Текст раздела...

\subsubsection{Проектирование модуля синхронизации текста и голоса}

Текст раздела...

\subsubsection{Разработка системы управления историей мультимодального диалога}

Текст раздела...

\subsubsection{Функциональная спецификация обработчика сообщений}

Текст раздела...

\subsubsection{Проектирование модуля потоковой передачи голоса}

Текст раздела...

\subsection{Проектирование модуля базы данных и интерфейсов}

\subsubsection{Архитектура взаимодействия с PostgreSQL}

Текст раздела...

\subsubsection{Проектирование выполнения запросов к уникальной базе данных}

Текст раздела...

\subsubsection{Разработка программного интерфейса для управления агентами}

Текст раздела...

\subsubsection{Функциональная спецификация мультимодального чат-интерфейса}

Текст раздела...

\subsubsection{Проектирование WebSocket-соединений для реального времени}

Текст раздела...

\subsection{Проектирование пользовательских интерфейсов}

\subsubsection{Архитектура веб-чата с поддержкой текста и голоса}

Текст раздела...

\subsubsection{Проектирование голосового интерфейса с визуальной обратной связью}

Текст раздела...

\subsubsection{Разработка шаблонов мультимодального пользовательского интерфейса}

Текст раздела...

\subsubsection{Функциональная спецификация голосовых команд}

Текст раздела...

\subsubsection{Проектирование обработчика мультимодального ввода}

Текст раздела...

\subsection{Выводы по главе 2}

Текст раздела...

\newpage

% ==================== ГЛАВА 3 ====================
\section{Реализация, интеграция и оценка эффективности мультимодальной агентной системы для уникальной базы данных}

\subsection{Реализация ядра системы и базовых механизмов}

\subsubsection{Программная реализация базового класса агента}

Текст раздела...

\subsubsection{Разработка системы регистрации и управления агентами}

Текст раздела...

\subsubsection{Реализация механизмов коммуникации}

Текст раздела...

\subsubsection{Программирование параллельного выполнения задач}

Текст раздела...

\subsection{Реализация аналитических агентов}

\subsubsection{Интеграция библиотек химического анализа}

Текст раздела...

\subsubsection{Реализация моделей машинного обучения для оценки токсичности}

Текст раздела...

\subsubsection{Разработка алгоритмов поиска структурных аналогов}

Текст раздела...

\subsubsection{Программная реализация методов предсказания свойств}

Текст раздела...

\subsubsection{Реализация анализа межвещественных взаимодействий}

Текст раздела...

\subsubsection{Разработка модуля оценки эффективности}

Текст раздела...

\subsection{Реализация хронобиологических агентов для уникальной базы данных}

\subsubsection{Создание базы знаний на основе единственной в мире базы хронобиотиков}

Текст раздела...

\subsubsection{Разработка алгоритмов извлечения информации из научных баз по хронобиологии}

Текст раздела...

\subsubsection{Реализация механизмов анализа экспериментальных данных}

Текст раздела...

\subsubsection{Программирование комплексного анализа веществ}

Текст раздела...

\subsubsection{Реализация интеллектуального поиска по уникальной базе данных}

Текст раздела...

\subsubsection{Разработка алгоритмов исследования механизмов}

Текст раздела...

\subsection{Реализация мультимодальных диалоговых агентов}

\subsubsection{Разработка мультимодальной диалоговой системы с поддержкой контекста}

Текст раздела...

\subsubsection{Реализация алгоритмов объяснения результатов в текстовом и голосовом форматах}

Текст раздела...

\subsubsection{Программирование вопросно-ответной системы с поддержкой голосового ввода}

Текст раздела...

\subsubsection{Разработка методов суммирования текстов с голосовым озвучиванием}

Текст раздела...

\subsubsection{Реализация рекомендательной системы с мультимодальным выводом}

Текст раздела...

\subsection{Реализация агентов поиска и мониторинга научной литературы}

\subsubsection{Интеграция с API PubMed для поиска научных статей}

Текст раздела...

\subsubsection{Разработка алгоритмов поиска в Google Scholar и CrossRef}

Текст раздела...

\subsubsection{Реализация системы мониторинга новых публикаций по теме}

Текст раздела...

\subsubsection{Программирование извлечения ключевой информации из найденных статей}

Текст раздела...

\subsubsection{Разработка механизмов обновления базы знаний на основе новых данных}

Текст раздела...

\subsubsection{Реализация уведомлений о появлении релевантных публикаций}

Текст раздела...

\subsection{Реализация исследовательских агентов}

\subsubsection{Разработка алгоритмов поиска экспериментальных исследований}

Текст раздела...

\subsubsection{Реализация методов генерации гипотез}

Текст раздела...

\subsubsection{Программирование систематического обзора литературы}

Текст раздела...

\subsubsection{Разработка исследования механизмов}

Текст раздела...

\subsubsection{Реализация отслеживания научных трендов в области хронобиологии}

Текст раздела...

\subsection{Реализация голосовых агентов для мультимодального чата}

\subsubsection{Интеграция модели Whisper для распознавания речи}

Текст раздела...

\subsubsection{Настройка и интеграция Text-to-Speech-движков для синтеза речи}

Текст раздела...

\subsubsection{Разработка алгоритмов обработки голосовых команд}

Текст раздела...

\subsubsection{Реализация обработки аудиопотоков в реальном времени}

Текст раздела...

\subsubsection{Интеграция мультимодального взаимодействия в единый интерфейс}

Текст раздела...

\subsubsection{Реализация синхронизации текстового и голосового вывода}

Текст раздела...

\subsection{Реализация многоязычных агентов}

\subsubsection{Интеграция библиотек определения языка}

Текст раздела...

\subsubsection{Реализация механизмов перевода для текста и голоса}

Текст раздела...

\subsubsection{Разработка системы локализации интерфейса}

Текст раздела...

\subsubsection{Реализация многоязычного диалога с поддержкой голоса}

Текст раздела...

\subsection{Реализация цитационных агентов и агентов данных}

\subsubsection{Разработка системы построения библиографии}

Текст раздела...

\subsubsection{Реализация извлечения цитат}

Текст раздела...

\subsubsection{Программирование форматирования ссылок}

Текст раздела...

\subsubsection{Разработка отслеживания источников}

Текст раздела...

\subsubsection{Реализация отслеживания контента}

Текст раздела...

\subsubsection{Разработка веб-скрининга}

Текст раздела...

\subsection{Реализация модуля химического анализа}

\subsubsection{Реализация расчета молекулярных дескрипторов}

Текст раздела...

\subsubsection{Разработка предсказания свойств веществ}

Текст раздела...

\subsubsection{Интеграция с внешними химическими базами данных для обогащения уникальной коллекции}

Текст раздела...

\subsubsection{Реализация распознавания химических структур из изображений}

Текст раздела...

\subsubsection{Разработка парсеров химических форматов}

Текст раздела...

\subsubsection{Поддержка голосового ввода химических названий}

Текст раздела...

\subsection{Реализация модуля графов знаний}

\subsubsection{Разработка извлечения сущностей из текстов и базы данных}

Текст раздела...

\subsubsection{Реализация построения графа знаний на основе реляционных данных хронобиотиков}

Текст раздела...

\subsubsection{Программирование гибридного поиска с учетом ограниченного объема данных}

Текст раздела...

\subsubsection{Разработка механизмов логического вывода для выявления новых взаимосвязей}

Текст раздела...

\subsubsection{Реализация визуализации графа с голосовыми подсказками}

Текст раздела...

\subsection{Реализация модуля генерации с дополненной выборкой}

\subsubsection{Разработка стратегий сегментации текстов}

Текст раздела...

\subsubsection{Реализация генерации векторных представлений}

Текст раздела...

\subsubsection{Программирование гибридного поиска}

Текст раздела...

\subsubsection{Разработка ранжирования результатов}

Текст раздела...

\subsubsection{Интеграция с векторными хранилищами}

Текст раздела...

\subsection{Интеграция большой языковой модели}

\subsubsection{Реализация унифицированного API для взаимодействия с различными LLM}

Текст раздела...

\subsubsection{Настройка и развертывание открытых моделей (BLOOM, LLaMA, Mistral, Qwen)}

Текст раздела...

\subsubsection{Интеграция с проприетарными моделями (OpenAI GPT, Anthropic Claude, Google Gemini)}

Текст раздела...

\subsubsection{Разработка системы роутинга запросов для выбора оптимальной модели}

Текст раздела...

\subsubsection{Реализация механизмов кэширования и оптимизации запросов}

Текст раздела...

\subsubsection{Создание библиотеки промпт-шаблонов для различных типов агентов}

Текст раздела...

\subsubsection{Дообучение моделей на предметных данных хронобиологии}

Текст раздела...

\subsubsection{Реализация инструментов расширения функциональности (LLM Tools)}

Текст раздела...

\subsection{Разработка гибридной системы поиска на основе уникальной реляционной базы}

\subsubsection{Создание векторных представлений данных из PostgreSQL}

Текст раздела...

\subsubsection{Реализация RAG-системы поверх реляционной базы данных}

Текст раздела...

\subsubsection{Построение графа знаний из реляционных данных}

Текст раздела...

\subsubsection{Разработка гибридного поискового механизма для эффективного извлечения знаний}

Текст раздела...

\subsubsection{Интеграция результатов поиска по литературе с данными из базы}

Текст раздела...

\subsection{Реализация модуля мультимодального чата}

\subsubsection{Разработка чат-движка с поддержкой текста и голоса}

Текст раздела...

\subsubsection{Реализация синхронизации различных модальностей}

Текст раздела...

\subsubsection{Программирование системы управления историей диалога}

Текст раздела...

\subsubsection{Разработка обработчика мультимодальных сообщений}

Текст раздела...

\subsubsection{Реализация потоковой передачи голоса}

Текст раздела...

\subsection{Реализация базы данных и программных интерфейсов}

\subsubsection{Разработка моделей данных в соответствии с ER-диаграммой}

Текст раздела...

\subsubsection{Реализация менеджера базы данных}

Текст раздела...

\subsubsection{Создание REST API для управления агентами}

Текст раздела...

\subsubsection{Разработка WebSocket-соединений для реального времени}

Текст раздела...

\subsubsection{Реализация миграций базы данных}

Текст раздела...

\subsection{Реализация пользовательских интерфейсов}

\subsubsection{Разработка веб-интерфейса с поддержкой текстового ввода/вывода}

Текст раздела...

\subsubsection{Создание голосового интерфейса с визуализацией}

Текст раздела...

\subsubsection{Реализация многоязычного интерфейса}

Текст раздела...

\subsubsection{Разработка визуализации графа знаний с интерактивными элементами}

Текст раздела...

\subsubsection{Создание адаптивного дизайна для мобильных устройств}

Текст раздела...

\subsubsection{Интеграция кнопок и индикаторов для голосового управления}

Текст раздела...

\subsubsection{Реализация панели мониторинга новых публикаций}

Текст раздела...

\subsection{Тестирование системы}

\subsubsection{Функциональное тестирование аналитических агентов}

Текст раздела...

\subsubsection{Тестирование агентов поиска научной литературы}

Текст раздела...

\subsubsection{Тестирование голосовых агентов на различных акцентах и языках}

Текст раздела...

\subsubsection{Тестирование многоязычных агентов на переводе научных терминов}

Текст раздела...

\subsubsection{Тестирование мультимодального взаимодействия}

Текст раздела...

\subsubsection{Интеграционное тестирование взаимодействия агентов}

Текст раздела...

\subsubsection{Нагрузочное тестирование}

Текст раздела...

\subsubsection{Сравнительное тестирование различных LLM на предметных задачах}

Текст раздела...

\subsection{Оценка эффективности мультимодальной системы для уникальной базы данных}

\subsubsection{Оценка точности химического анализа на ограниченном наборе данных}

Текст раздела...

\subsubsection{Оценка эффективности поиска новых научных публикаций}

Текст раздела...

\subsubsection{Оценка качества распознавания и синтеза речи}

Текст раздела...

\subsubsection{Оценка точности перевода научной терминологии}

Текст раздела...

\subsubsection{Оценка релевантности результатов гибридного поиска в уникальной базе}

Текст раздела...

\subsubsection{Сравнительный анализ качества ответов различных LLM}

Текст раздела...

\subsubsection{Оценка пользовательского опыта при мультимодальном взаимодействии}

Текст раздела...

\subsubsection{Сравнительный анализ с альтернативными подходами доступа к данным}

Текст раздела...

\subsection{Анализ результатов внедрения уникальной системы}

\subsubsection{Статистика использования системы на реальных данных}

Текст раздела...

\subsubsection{Обратная связь от исследователей-хронобиологов}

Текст раздела...

\subsubsection{Анализ эффективности обнаружения новых научных публикаций}

Текст раздела...

\subsubsection{Анализ эффективности различных LLM в контексте предметной области}

Текст раздела...

\subsubsection{Выявленные ограничения и пути их преодоления}

Текст раздела...

\subsubsection{Научная значимость созданной системы как единственного интеллектуального интерфейса к уникальной базе данных}

Текст раздела...

\subsection{Перспективы развития системы}

\subsubsection{Расширение набора агентов для смежных областей}

Текст раздела...

\subsubsection{Интеграция с дополнительными внешними источниками для обогащения данных}

Текст раздела...

\subsubsection{Улучшение алгоритмов поиска и анализа научной литературы}

Текст раздела...

\subsubsection{Поддержка новых LLM и мультимодальных моделей}

Текст раздела...

\subsubsection{Развитие мультимодальных интерфейсов}

Текст раздела...

\subsubsection{Улучшение качества синтеза и распознавания речи}

Текст раздела...

\subsubsection{Планы по развитию и пополнению уникальной базы данных хронобиотиков}

Текст раздела...

\subsection{Выводы по третьей главе}

Текст раздела...

\newpage

% ==================== ЗАКЛЮЧЕНИЕ ====================
\section*{Заключения}
\addcontentsline{toc}{section}{Заключения}



\newpage

% ==================== БИБЛИОГРАФИЧЕСКИЙ СПИСОК ====================
\section*{Библиографический список}
\addcontentsline{toc}{section}{Библиографический список}

\begin{thebibliography}{99}
	

\end{thebibliography}

\newpage

% ==================== ПРИЛОЖЕНИЕ А ====================
\section*{Приложение А. Архитектурные диаграммы мультимодальной агентной системы}
\addcontentsline{toc}{section}{Приложение А. Архитектурные диаграммы мультимодальной агентной системы}

Текст приложения А...

\newpage

% ==================== ПРИЛОЖЕНИЕ Б ====================
\section*{Приложение Б. Детальное описание структуры уникальной базы данных хронобиотиков}
\addcontentsline{toc}{section}{Приложение Б. Детальное описание структуры уникальной базы данных хронобиотиков}

Текст приложения Б...

\end{document}